%% ============================================================
%% 分片 20：原文第 7 节「Fluxes and $S$-Duality」（原文 glshort.tex 起于第 5685 行）
%% 由 tools/harvmac2tex.py 从 arXiv 源码（hep-th/0604151v3）自动生成。
%% 只译英文散文；公式、编号、\label、\cite 一律不改。
%% ============================================================

\subsection{Compactification To Two Dimensions} \label{sec:comptwop}

Now we compactify to two dimensions. Consider ${\EUN}=4$ super
Yang-Mills theory on $M=\Sigma\times C$ with gauge group
$G_{\rm{ad}}$. The data in this theory are a $G_{\rm{ad}}$-bundle
$E\to\Sigma\times C$, a connnection $A$ on $E$, and various other
fields.  Any choice of $E$ determines a class $\xi(E)\in
H^2(\Sigma\times C,\CZ)$.  Since $H^2(\Sigma\times
C,\CZ)=H^2(\Sigma,\CZ)\oplus H^1(\Sigma,H^1(C,\CZ))\oplus
H^2(C,\CZ),$ we have % [WARN] 公式 \eqn\eqref{eq:ytro} 无编号或重复
\begin{equation*}
\tag*{(7.18)}
\xi(E)\in H^2(\Sigma,\CZ)\oplus
H^1(\Sigma,H^1(C,\CZ))\oplus H^2(C,\CZ).
\end{equation*}
 Relative to this
decomposition we write
$\xi(E)=\xi(E)^{2,0}+\xi(E)^{1,1}+\xi(E)^{0,2}$.

We can describe the same thing in the low energy effective
sigma-model.  Here we consider low energy data on $\Sigma\times
C$, which when restricted to $p\times C$ for a point $p\in \Sigma$
determine a point in $\MHAD$.  As $p$ varies, we get a map
$\Phi:\Sigma\to \MHAD$.   Composing $\Phi$ with the identity map
on $C$, we get a map $\Phi\times 1:\Sigma\times C\to \MHAD\times
C$. Now recall the universal bundle ${\EUE}_{\rm ad}$ and the
class $\zeta=\xi({\EUE}_{\rm {ad}})\in H^2( \MHAD\times C,\CZ)$.
Pulling it back to $\Sigma\times C$, we get a class $(\Phi\times
1)^*(\zeta)\in H^2(\Sigma\times C,\CZ)$. By chasing through the
definitions, one can see that % 式 (7.19)
\setcounter{equation}{18}
\begin{equation}\label{eq:yrog}
\xi(E)=(\Phi\times
1)^*(\zeta).
\end{equation}
 In fact, the universal bundle ${\EUE}_{ad}$ is
defined so that $(\Phi\times 1)^*({\EUE}_{{\rm ad}})=E$, so
$(\Phi\times 1)^*(\zeta)=(\Phi\times 1)^*(\xi(\EUE_{{\rm ad}}))
=\xi((\Phi\times 1)^*(\EUE_{{\rm ad}}))=\xi(E)$.

Let us further specialize to $\Sigma=S^1\times \tilde S^1$, where
the two factors are circles. We write $W=\tilde S^1\times C$. So
$\Sigma\times C=S^1\times W$. $\xi(E)$ has the familiar expansion
$\xi(E)={\bf a}+{\bf m}$ with ${\bf a}\in H^1(W,\CZ)$, ${\bf m}\in
H^2(W,\CZ)$. Now, however, we can further expand
$H^1(W,\CZ)=H^1(\tilde S^1,\CZ)\oplus H^1(C,\CZ)$, so we have
${\bf a}={\bf a}_0+{\bf a}_1$ with % 式 (7.20)
\setcounter{equation}{19}
\begin{equation}\label{eq:mimot}
\begin{aligned}{\bf
a}_0&\in H^1(\tilde S^1,\CZ) =\CZ,\\ {\bf a}_1& \in
H^1(C,\CZ).\\\end{aligned}
\end{equation}
 And similarly, we can expand
$H^2(W,\CZ)=H^1(\tilde S^1,H^1(C,\CZ))\oplus H^2(C,\CZ)$, so we
have ${\bf m}={\bf m}_0+{\bf m}_1$ with % 式 (7.21)
\setcounter{equation}{20}
\begin{equation}\label{eq:nimot}
\begin{aligned}{\bf
m}_0&\in H^2(C,\CZ)=\CZ \\
                     {\bf m}_1&\in
                     H^1(\tilde S^1,H^1(C,\CZ))=H^1(C,\CZ).\\\end{aligned}
\end{equation}

Comparing back to the expansion \eqref{eq:ytro}, we see that
% 式 (7.22)
\setcounter{equation}{21}
\begin{equation}\label{eq:zimot}
\xi(E)^{2,0}={\bf a}_0,~\xi(E)^{0,2}={\bf m}_0.
\end{equation}

We can similarly decompose the character $\bf e$ that is dual to
$\bf a$. Clearly, we have ${\bf e}={\bf e}_0+{\bf e}_1$, where
${\bf e}_0$ is a character of ${\bf a}_0\in \CZ$ and ${\bf e}_1$
is a character of ${\bf a}_1\in H^1(C,\CZ)$.  Hence, using
 the selfduality of $\CZ$ and Poincar\'e duality, % 式 (7.23)
\setcounter{equation}{22}
\begin{equation}\label{eq:tew}
\begin{aligned} {\bf e}_0 & \in
\CZ^\vee=\CZ\\ {\bf e}_1&\in H^1(C,\CZ)^\vee=H^1(C,\CZ).\\\end{aligned}
\end{equation}
 In
view of \eqref{eq:coco}, the transformation under $S$-duality is
% 式 (7.24)
\setcounter{equation}{23}
\begin{equation}\label{eq:niko}
\begin{aligned}({\bf e}_0,{\bf m}_0)&\to (-{\bf m}_0,{\bf
e}_0)\\ ({\bf e}_1,{\bf m}_1)&\to (-{\bf m}_1,{\bf e}_1).\\\end{aligned}
\end{equation}

Our next goal is to interpret ${\bf e}_0$, ${\bf e}_1$, ${\bf
m}_0$, and ${\bf m}_1$ in the effective two-dimensional
sigma-model with target $\MHAD$.  We will consider both the case
of a closed Riemann surface $\Sigma$, and the case when $\Sigma$
has a boundary. The latter case will enable us to understand the
implications of \eqref{eq:niko}\ for the geometric Langlands program.

\par\medskip{\it Interpretation Of ${\bf m}_0$}

The easiest to interpret is ${\bf m}_0=\xi(E)^{0,2}$.  The target
space $\MH(G_{\rm ad},C)$ of our sigma-model is not connected. Its
components are labeled by the topological type of the $G_{\rm
{ad}}$-bundle $E\to C$.  But this is exactly what is measured by
${\bf m}_0$.

\par\medskip{\it Interpretation Of ${\bf e}_0$}

Consider in general a sigma-model of maps $\Phi:\Sigma\to X$, for
some $X$.  A flat $B$-field is an element  $b\in H^2(X,U(1))$. A
flat $B$-field is incorporated in the sigma-model path integral as
follows, for the case that $\Sigma$ has no boundary.  Given a map
$\Phi:\Sigma\to X$, one pulls back $b$ to $\Phi^*(b)\in
H^2(\Sigma,U(1))=U(1)$, and then one includes in the path integral
a factor of $\Phi^*(b)$.  Thus, in this situation, incorporating
the flat $B$-field has the effect of weighting by phases the
different components of maps of $\Sigma$ to $X$.

Now ${\bf e}_0$ determines a flat $B$-field in the sigma-model of
maps $\Phi:\Sigma\to \MH(G_{\rm ad},C)$. Indeed, there is as we
have explained a natural class $\zeta\in H^2(\MH(G_{\rm
ad},C),\CZ)$, which expresses the  obstruction to lifting the
universal $G_{\rm{ad}}$ bundle  ${\EUE}_{\rm {ad}}\to\MH(G_{\rm
ad},C)\times C$ to a $\bar G$-bundle.  By composing $\zeta$ with ${\bf
e}_0:\CZ\to U(1)$, we get a flat $B$-field $b_{{\bf e}_0}={\bf
e}_0(\zeta)\in H^2(\MH(G_{\rm{ad}},C),U(1))$.  We claim that the
role of ${\bf e}_0$ in the effective sigma-model of maps
$\Sigma\to \MHAD$ is precisely to weight every map in the way that
one would expect for a sigma-model with the flat $B$-field
$b_{{\bf e}_0}$.

In fact, the definition of ${\bf e}_0$ is that the path integral
includes a phase factor ${\bf e}_0({\bf a}_0)={\bf
e}_0(\xi(E)^{2,0})$.   But $\xi(E)^{2,0}=\Phi^*(\zeta)$; this
follows upon restricting \eqref{eq:yrog}\ to a point in $C$.
So the phase factor induced in the path integral by
${\bf e}_0$ is ${\bf e}_0(\Phi^*(\zeta))=\Phi^*({\bf
e}_0(\zeta))=\Phi^*(b_{{\bf e}_0})$.  This justifies  our claim
that the effect of ${\bf e}_0$ in the sigma-model is to generate a
flat $B$-field $b_{{\bf e}_0}$.

\eqref{eq:niko}\ therefore means that $S$-duality exchanges the topological
class ${\bf m}_0$ of a flat $G_{\rm{ad}}$-bundle with the flat
$B$-field determined by ${\bf e}_0$.

\par\medskip{\it Incorporation Of Branes}

For applications to the geometric Langlands program, the real
payoff is to understand the implications of all this for branes.

So we take $\Sigma$ to be a Riemann surface with boundary.
$M=\Sigma\times C$ is  therefore a four-manifold with boundary
$\partial M=\partial\Sigma\times C$.  On $\partial M$, we place
some supersymmetric boundary condition.

The effective two-dimensional description is a sigma-model of maps
$\Phi:\Sigma\to \MHAD$, with boundary condition corresponding to
some brane.  We would like to understand the role of ${\bf e}_0$
and ${\bf m}_0$ in this description.

There is little new to say about ${\bf m}_0$.  It labels the
components of $\MHAD$, whether $\Sigma$ has a boundary or not.

The role of ${\bf e}_0$ is more subtle.  As we have seen, the
sigma-model with target $\MHAD$ is endowed with a flat $B$-field
${\bf e}_0(\zeta)$. A flat $B$-field has a very interesting effect
on branes \cite{ref102}.  In the absence of a
$B$-field, a brane on $\MHAD$ has a Chan-Paton bundle, which is a
vector bundle over $\MHAD$  (or more generally a sheaf, perhaps
supported on a submanifold, that defines a $K$-theory class of
$\MHAD$). However, in the presence of a flat $B$-field, associated
with an element $b\in H^2(\MHAD,U(1))$, the Chan-Paton bundle
becomes a twisted vector bundle (or more generally a twisted sheaf
related to an element of the twisted $K$-theory of $\MHAD$),
twisted by $b$ in the sense of eqn. \eqref{eq:hofto}.  It is because of this
that we introduced the concept of a twisted vector bundle.

So in short, for ${\bf e}_0\not=0$,  the Chan-Paton bundle of a
brane is a twisted vector bundle, twisted by  $b_{{\bf e}_0}={\bf
e}_0(\zeta)$.  Luckily, from the analysis at the end of section
第~\ref{sec:reviewthis}~节, we have a plentiful supply of such twisted vector
bundles.  If $\varrho$ is any irreducible representation of $\bar
G$ such that the character of the center of $\bar G$ defined by
$\varrho$ is equal to ${\bf e}_0$, then the universal bundle
${\EUE}_\varrho$ in the representation $\varrho$ is an example of
a twisted vector bundle for the flat $B$-field $b_{{\bf e}_0}$.

Other important examples of twisted branes can be constructed by
picking  a submanifold  $Y\subset \MHAD$ such that $b_{{\bf e}_0}$
is trivial when restricted to $Y$.  In this case, the definition
of a brane supported on $Y$ is independent of ${\bf e}_0$ (up to a
not quite canonical isomorphism).  For example, $Y$ could be a
point in $\MHAD$. Certainly $b_{{\bf e}_0}$ is trivial when
restricted to a point, so zerobranes exist for any ${\bf e}_0$.
{}From such zerobranes, we can form electric eigenbranes, as we
explain in section 第~\ref{sec:electricbranes}~节. A slightly more subtle example
is a magnetic eigenbrane, a brane of rank one supported on a fiber
${\CMF}$ of the Hitchin fibration. The Chan-Paton bundle of such a
brane should be a flat\footnote{One can here interpret flatness to
mean, just as for ordinary line bundles, that the transition
functions are constants.} twisted line bundle. In \cite{ref21}, it is
shown that the flat $b$-field $b_{{\bf e}_0}$ is trivial when
restricted to ${\CMF}$, as a result of which the space of flat
twisted line bundles is isomorphic, but not canonically
isomorphic, to the space of ordinary flat line bundles on
${\CMF}$.

We can now deduce from $S$-duality a statement about branes. The
duality transformation $S$ maps $({\bf e}_0,{\bf m}_0)\to (-{\bf
m}_0,{\bf e}_0)$, so it exchanges the topology of the component of
$\MH$ on which a brane is supported with the flat $B$-field
$b_{{\bf e}_0}$ by which its Chan-Paton bundle is twisted. It
also, of course, exchanges $G_{\rm ad}$ with\footnote{By
$^L\negthinspace G_{\rm ad}$, we mean the adjoint form of the
group $^L\negthinspace G$. The statement we are describing here is
best expressed in terms of adjoint bundles on both sides to allow
all possible topologies. Momentarily we indicate explicitly
whether we want a given component of $\MH$ or its universal
cover.} $^L\negthinspace G_{{\rm ad}}$, and (as we explain in the
concluding remark of this section), exchanges $\MH(G_{{\rm ad}})$
with $\hat\M_H({}^L\negthinspace G_{{\rm ad}})$, which we define
to be the universal cover of $\MH({}^L\negthinspace G_{{\rm
ad}})$.

As a special case of this duality, a point on one side, contained
in a fiber $\CMF$ of the Hitchin fibration, is mapped on the other
side to a brane of rank one supported on the corresponding fiber
$^L\negthinspace\CMF=\Xi(\CMF)$ of the Hitchin fibration, and
endowed with a flat twisted line bundle. $\CMF$ is a union of
complex tori, labeled by the characteristic class ${\bf
m}_0=\xi(E)$ of the Higgs bundle. The choice of a component of
$\CMF$ on one side determines on the other side  the discrete
electric field ${\bf e}_0$ and hence the twist. The choice of a
point on $\CMF$ determines a flat twisted line bundle on
$^L\negthinspace\CMF$ (to which it maps under the duality
transformation $S$). Of course, this relationship between $\CMF$
and $^L\negthinspace\CMF$ is reciprocal.  This twisted duality
between $\CMF$ and $^L\negthinspace\CMF$ is in fact one of the
main results of Hausel and Thaddeus \cite{ref21}.

\par\medskip{\it Interpretation Of ${\bf e}_1$ And ${\bf
m}_1$}

\def\MHBARG{\MH(\bar G)}

Finally, let us discuss the interpretation of ${\bf e}_1$ and
${\bf m}_1$ in two-dimensional terms. It will be helpful to begin
by comparing the theories with gauge groups $\bar G$ and
$G_{\rm{ad}}$.

$\MHBARG$ is simply-connected.  It has a natural group of
symmetries ${\bf E}_C=H^1(C,\CZ)$.  Indeed, ${\bf E}_C$
parametrizes $\CZ$-bundles over $C$.  A $\bar G$ Higgs bundle
$(E,\varphi)$ can be tensored with a $\CZ$-bundle to make a new
$\bar G$ Higgs bundle.  (Concretely, this operation multiplies the
holonomies of $E$ around noncontractible loops in $C$ by elements
of the center of $\bar G$.)  So this gives a group ${\bf E}_C$ of
symmetries of the sigma-model with target $\MHBARG$. The Hilbert
space of this sigma-model can be decomposed in characters of ${\bf
E}_C$.

On the other hand, $\MHAD$ has no such geometrical symmetries. But
it has a fundamental group ${\bf M}_C=H^1(C,\CZ)$, which is
isomorphic to ${\bf E}_C$. To understand where this fundamental
group comes from, a shortcut is to note that one component
$\MHAD_0$ of $\MHAD$, namely the component that parametrizes Higgs
bundles that can be lifted to $\bar G$, is simply
$\MHAD_0=\MHBARG/{\bf E}_C$.  Dividing by ${\bf E}_C$ eliminates
the geometrical symmetries of $\MHBARG$, but of course it creates
a fundamental group.\footnote{ ${\bf E}_C$ does not act freely, so $\M_H(G_{ad})$
 has orbifold singularities.  (It also has more severe singularities from reducible
 Higgs bundles.)
As is familiar in
sigma models, the fundamental group of $\M_H(G_{ad})$ must be understood in
an orbifold
sense.  Alternatively, one can rely on the four-dimensional gauge theory
instead of reducing
to the sigma model with its singularities.} So $\pi_1(\MHAD_0)\cong {\bf E}_C$.

Actually, the fundamental group is the same for any component of
$\MHAD$.  $\MHAD$ is defined by dividing the space of all
solutions of Hitchin's equations, for gauge group $G_{\rm {ad}}$,
by the group ${\cal G}_{\rm{ad}}(C)$ of all $G_{\rm{ad}}$-valued
gauge transformations on $C$. If one were to divide only by the
connected component of ${\cal G}_{\rm{ad}}(C)$, one would get the
universal cover $\hat\M_H(G_{\rm {ad}},C)$.  The fundamental group
of $\MHAD$ is therefore the group of components of ${\cal G}_{\rm
ad}(C)$, and this is ${\bf E}_C=H^1(C,\CZ)$, for any component of
$\MHAD$.

In sum, strings moving on $\MHAD$ have a discrete group of
conserved winding numbers ${\bf M}_C=\pi_1(\MHAD)=H^1(C,\CZ)$.
Likewise strings moving on $\MHBARG$ have a discrete group of
conserved momenta ${\bf E}_C=H^1(C,\CZ)$.

Let us compare the symmetries ${\bf E}_C$ and ${\bf M}_C$ to what
we can see in the underlying gauge theory. In $G_{\rm{ad}}$ gauge
theory on $M=\Sigma\times C=S^1\times \tilde S^1\times C$, the
bundle $E$ and other data determine a map $\Phi:S^1\times \tilde
S^1\to \MHAD$. In analyzing the topology of this situation, we
expanded the characteristic class $\xi(E)$ as $\xi(E)={\bf
a}\oplus {\bf m}$, where ${\bf m}$ is the restriction of $\xi(E)$
to $\tilde S^1\times C$ (more precisely, to $p\times \tilde
S^1\times C$, for a point $p\in S^1$). In the low energy
sigma-model, ${\bf a}$ and ${\bf m}$ are invariants of
$\Phi:\Sigma\to \MHAD$. In particular, ${\bf m}$ only depends on
the restriction $\Phi|_{\tilde S^1}$ of $\Phi$ to $\tilde S^1$.

In \eqref{eq:nimot}, we further expanded ${\bf m}={\bf m}_0\oplus {\bf m}_1$
with ${\bf m}_1\in H^1(\tilde S^1,H^1(C,\CZ))$.  ${\bf m}_1$ is a
topological invariant of $\Phi|_{\tilde S^1}$ that  vanishes for
constant maps of $\tilde S^1$ to $\MHAD$. So it measures the
homotopy class of the map $\Phi|_{\tilde S^1}$ in
$\pi_1(\MHAD)=H^1(C,\CZ)$.

Similarly, we could exchange the role of $S^1$ and $\tilde S^1$.
The restriction of $\xi(E)$ to $S^1 \times C$ (that is, to
$S^1\times q\times C$, for a point $q\in \tilde S^1$), is in the
above notation ${\bf m}_0\oplus {\bf a}_1$. In particular, ${\bf
a}_1\in H^1(S^1,H^1(C,\CZ))$ measures the winding of $\Phi$ in the
$S^1$ or ``time'' direction. The character ${\bf e}_1$ which is
dual to ${\bf a}_1$ therefore measures the conserved momentum of
the strings.

\par\medskip{\it An Example}

For an important illustration of all this, consider the Langlands
dual pair $G_{\rm{ad}}$ and $^L\negthinspace\bar G$.  In
$G_{\rm{ad}}$ gauge theory, we set ${\bf e}=0$ and consider a
sigma-model with target $\MHAD$.  In this model, there are no
conserved momenta, but there is a symmetry ${\bf M}_C$ of string
windings.

In the dual picture, the gauge group is $^L\negthinspace\bar G$,
we set ${\bf m}=0$, and the sigma-model with target $\MHBARG$ has
no conserved string windings, but a symmetry group ${\bf E}_C$ of
discrete conserved momenta.

$S$-duality or Montonen-Olive or Langlands duality exchanges the
two pictures, exchanging ${\bf E}_C({}^L\negthinspace\bar G)$ with
${\bf M}_C(G_{\rm{ad}})$. The fact that the duality exchanges the
discrete conserved momenta and windings of strings is an aspect of
its relation to $T$-duality in two dimensions.

\par\medskip{\it A Concluding Comment}

More generally, we can  simply specify ${\bf m}_0$ and ${\bf e}_0$
as we please. Then we consider branes on a component of $\MHAD$
labeled topologically by ${\bf m}_0$; the Chan-Paton bundles of
the branes are twisted by ${\bf e}_0$.

We still have two ways to proceed with the quantization.  If we
divide by all $G_{\rm {ad}}$-valued gauge transformations, then
the target space is the component of $\MH(G_{\rm ad},C)$ labeled
by ${\bf m}_0$.  In this case, there is a finite group ${\bf M}_1$
that classifies the string winding numbers, but there is no group
${\bf E}_1$ of geometrical symmetries.  If we divide by only the
connected gauge transformations, then the target is the universal
cover $\hat\M_H(G_{\rm ad},C)$.  In this case, there is a group
${\bf E}_1$ of geometrical symmetries, but no group ${\bf M}_1$ of
string winding numbers.  $S$-duality exchanges ${\bf E}_1$ and
${\bf M}_1$, so (in addition to exchanging $G$ and
$^L\negthinspace G$) it exchanges the two methods of quantization.

In particular, when we apply $S$-duality to branes of specified
${\bf e}_0$ and ${\bf m}_0$, we must exchange the two methods of
quantization in addition to exchanging the two adjoint groups
$G_{\rm ad}$ and $^L\negthinspace G_{\rm ad}$.
